\section{Validating an energetic epitope definition}\label{sec:energetic}

Before an energetic definition can drive design, its per-residue energies have to be trustworthy, so
we test them where the answer is already known. The pilot is 3HFM, the HyHEL-10 Fab bound to hen egg
lysozyme, for which SKEMPI records experimental alanine-scanning $\Delta\Delta G$ on thirteen antigen
residues. Three dominate --- Lys96 ($+6.5$ kcal/mol), Lys97 ($+5.6$), Tyr20 ($+4.3$) --- and the rest
sit at or below about 1 kcal/mol. We chose 3HFM because it is a well-measured hot-spot system and
because lysozyme is already a system with a validated molecular-dynamics setup to build on; none of
our own antigen structures appear in SKEMPI.

The method under test is a direct per-residue interaction energy. We run a 20 ns simulation of the
solvated complex and, for each antigen residue, average its nonbonded interaction with the antibody
(reaction-field Coulomb plus Lennard-Jones) over the trajectory. This is an interaction energy, not a
binding free energy: it omits the desolvation a residue pays to make a contact and any entropic
contribution. We therefore treat it as a ranking signal and ask one specific question of it.

The question is asymmetric, by design. For design we need to know which residues to fix by identity,
so a \emph{false positive} --- a residue the energy calls important but that is not --- is expensive:
we would spend design freedom locking a residue that does not carry the binding. A \emph{false
negative} is cheap: a residue that matters through structure, orientation, or entropy, which a
per-residue energy cannot see, is exactly the kind we hold by shape rather than identity, a
Category-2 residue. So we do not require the energy to recover every hot spot. We require that
everything it flags is real. The criterion is precision, not recall: no residue with high interaction
energy may have a low measured $\Delta\Delta G$.

The plan is an empirical fit. We tune the method --- the cutoff, which terms enter, whether the
solvent channel is included --- to empty the false-positive region on 3HFM, then test whether the same
method generalizes to other antibody--antigen complexes. If direct interaction energy alone cannot be
made free of false positives, that is the evidence that desolvation or water-mediated contributions
have to enter, and we move to a fuller treatment.

\subsection{What the interaction energy alone recovers}
The bare interaction energy recovers the strongest hot spots --- Lys97 and Lys96 are its top two
antigen calls, and Tyr20 and Asp101 also come out favorable --- but it produces a clear false
positive at Arg73, which ranks as the second-strongest interaction of all 129 residues yet has a
measured $\Delta\Delta G$ of only $-0.33$ kcal/mol (Figure~\ref{fig:ie_vs_ddg}). This is the failure a
bare interaction energy is expected to make: Coulomb credits a charged surface residue for its
contact without charging it the cost of stripping away the water it was hydrogen-bonded to in
solution. Removing the Coulomb term entirely and ranking on Lennard-Jones alone demotes Arg73 and
lifts the rank correlation with $\Delta\Delta G$ from $0.46$ to $0.73$, but it then over-calls
well-packed residues that are not hot spots (Leu75) and, at a fixed high-energy cut, multiplies the
false-positive count rather than reducing it. Electrostatics over-calls charged residues and packing
over-calls buried ones; both point at the same missing quantity, the net free-energy cost of
desolvation.

\subsection{Adding desolvation: per-residue MM-GBSA}
MM-GBSA is the principled fix. It keeps the Coulomb and Lennard-Jones interaction but adds the
generalized-Born polar-desolvation term, so a residue's raw contact is netted against the cost of
shedding its solvent; run single-trajectory, it decomposes onto the same 20 ns trajectory with no new
simulation. On the residue that motivated it, the mechanism works exactly as intended: Arg73's
favorable Coulomb contribution ($-58.7$ kcal/mol) is almost entirely cancelled by the $+57.2$ it pays
to desolvate, so its net per-residue contribution falls to $-3.6$ and it drops from the second-strongest
raw call to the flag threshold (Figure~\ref{fig:ie_vs_ddg_mmgbsa}).

The same trajectory sampled at twenty frames, however, is too sparse for the generalized-Born term to
be quantitatively reliable. It over-cancels the largest hot spot: Lys96's Coulomb ($-34.0$) is
overtaken by a $+38.9$ desolvation penalty, leaving it net \emph{unfavorable} ($+2.1$) --- a false
negative on precisely the residue we most need to recover. The rank correlation collapses to $0.05$,
and two residues (Arg73, Leu75) still sit marginally inside the false-positive zone. The desolvation
term is doing the right physics --- it is the only one of the three methods that nets a charged
residue's contact against its solvation --- but a single 20 ns replica sampled every nanosecond does
not give a per-residue generalized-Born estimate stable enough to lock a design against.

\subsection{Reading}
No method cleanly passes the precision test on 3HFM yet, and the three fail differently
(Table~\ref{tab:ie_methods}): the bare interaction energy over-calls charged residues, Lennard-Jones
alone over-calls packed ones, and desolvation-aware MM-GBSA corrects the charged over-call in
direction but, at this sampling, trades it for noise that sinks the strongest hot spot. Two follow-ups are now open: denser sampling of the existing MM-GBSA trajectory and a
comparison with a learned mutation-effect estimator. Denser sampling tests sensitivity to frame
selection; it cannot by itself distinguish systematic GB error from inadequate conformational
sampling. Neither approach is yet allowed to classify residues into the three categories. That is a
``find-what-works-on-this-system'' step, deliberately ahead of the question of what is cheap enough to
run at design scale.

\begin{table}[h]
\centering
\caption{Per-residue methods on 3HFM against thirteen measured antigen $\Delta\Delta G$. Spearman
$\rho$ is over the measured residues; false positives are measured residues above the $90^{\text{th}}$
interaction-energy percentile with $\Delta\Delta G < 1$ kcal/mol. Reproduce with
\texttt{energetics/plot\_ie\_vs\_ddg.py} at \texttt{-{}-channel} \texttt{ab\_total}, \texttt{ab\_lj},
\texttt{mmgbsa\_dg}.}
\label{tab:ie_methods}
\begin{tabular}{lccl}
\toprule
Method & Spearman $\rho$ & False positives & Characteristic failure \\
\midrule
Interaction energy (Coulomb + LJ) & $+0.46$ & 4 & over-calls charged residues (Arg73) \\
Lennard-Jones only                & $+0.73$ & 7 & over-calls packed residues (Leu75) \\
MM-GBSA (20 frames)               & $+0.05$ & 2 & demotes Arg73; over-desolvates Lys96 \\
\bottomrule
\end{tabular}
\end{table}

\begin{figure}[h]
\centering
\figorbox{ie_vs_ddg_3hfm.png}{0.9}
\caption{Bare per-residue interaction energy (Coulomb + LJ) against experimental $\Delta\Delta G$ for
3HFM. Lower panel: the false-positive test. The shaded lower-right region --- high interaction energy,
low measured $\Delta\Delta G$ --- must be empty; Arg73 sits inside it, the charged over-call the
desolvation step must fix. Generated by \texttt{energetics/plot\_ie\_vs\_ddg.py}
(\texttt{-{}-channel ab\_total}) from \texttt{holo\_ie.py} output.}
\label{fig:ie_vs_ddg}
\end{figure}

\begin{figure}[h]
\centering
\figorbox{ie_vs_ddg_3hfm_mmgbsa.png}{0.9}
\caption{Desolvation-inclusive per-residue MM-GBSA ($\Delta G$) against experimental $\Delta\Delta G$
for 3HFM, same trajectory, 20 frames. Arg73 is pulled back to the threshold (desolvation cancels its
Coulomb), but Lys96 --- the largest measured hot spot --- is over-desolvated to a net-unfavorable
$\Delta G$ and strands in the upper-left, and the rank correlation collapses ($\rho = 0.05$): the
mechanism is right but 20-frame single-replica GB is too noisy. Generated by
\texttt{energetics/plot\_ie\_vs\_ddg.py} (\texttt{-{}-channel mmgbsa\_dg}) from
\texttt{mmgbsa\_decomp\_to\_csv.py} output.}
\label{fig:ie_vs_ddg_mmgbsa}
\end{figure}


\subsection{Learned binding-$\Delta\Delta G$: a frozen Tamarind pilot}
\label{sec:tamarind_pilot}
\paragraph{Question and scope.}
We are testing whether an inexpensive learned predictor of mutation-induced binding free-energy
change can provide useful residue-importance information alongside the simulation-based methods.
The initial experiment is deliberately small: one StaB-ddG job evaluates two separate alanine
substitutions in the antigen of 3HFM. Lys96 is the strongest measured hotspot in our stored antigen
scan and is missed by the current MM-GBSA decomposition; Arg73 is the conspicuous raw-interaction
energy false positive. Selecting these two known disagreements makes this a diagnostic of specific
failure modes, not an unbiased estimate of model performance. The model is not being adopted as a
design standard on the basis of this run.

\paragraph{Experimental reference and observable.}
The stored SKEMPI-derived reference defines
\begin{equation}
\Delta\Delta G_{\mathrm{bind}}
=RT\ln\!\left(K_{d,\mathrm{mut}}/K_{d,\mathrm{WT}}\right),
\end{equation}
so positive values indicate weaker binding after mutation. The exact entries are K96A,
$+6.49$ kcal/mol, and R73A, $-0.33$ kcal/mol, each represented by one measurement in
\path{energetics/skempi_3hfm_ddg.csv}. The command below reads these entries and verifies
the structure checksum; the values are not inferred from a plotted point.
A predicted mutational $\Delta\Delta G$ is a different observable from a wild-type residue's
interaction energy or MM-GBSA decomposition. We will therefore compare mutation effects directly
with experiment, and compare simulation-derived signals as residue rankings; we will not subtract
these unlike quantities to obtain an error in kcal/mol. Model output units and sign convention
must be verified before any numerical comparison. Nor does success establish that the predictor
isolates an enthalpic contribution: mutation effects can include packing, solvation, and ensemble
changes.

\paragraph{Frozen structure and mutation mapping.}
The input is the existing RCSB 3HFM deposition at
\texttt{energetics/md/3hfm/structures/3hfm.pdb}, submitted as its full inline PDB text.
No local minimization, side-chain replacement, or trajectory-frame selection is performed for this
run. The interface partners are \texttt{binder1Chains=[H,L]} (HyHEL-10 Fab) and
\texttt{binder2Chains=[Y]} (lysozyme). The preparation script checks that Y:96 is LYS and
Y:73 is ARG in the coordinate records. The mutation array is
\texttt{["KY96A","RY73A"]}: two separate single mutants, not a combined double mutant.
The service's authenticated schema states that chain numbering is automatically converted to
start at residue 1 and that the remapped mutation is returned. The original and returned mutation
strings must be checked together before joining predictions to experiment.
The full SHA-256 of the input bytes is
\begin{verbatim}
92a6153cd4970117914d0767ee52838c860b1c9e6a05833b3f390bf502a2b3e1
\end{verbatim}
The same input text is frozen in \texttt{request.json}; \texttt{input\_manifest.json} records
its source path and checksum.

\paragraph{Service settings and execution record.}
The exact API tool identifier is \texttt{stabddg}. We explicitly set
\texttt{mcSamples=20}, the documented service default for Monte Carlo averaging, and
\texttt{seed=42}, a reproducibility setting rather than a tuned scientific parameter.
We retrieved and saved the authenticated settings schema from
\url{https://app.tamarind.bio/api/tools/stabddg/schema}, then submitted the request to the free
\texttt{validate-job} endpoint. Validation returned \texttt{valid=true}; the normalized settings
were saved and passed to \texttt{submit-job}. The job name is
\begin{verbatim}
episcaf-v3-3hfm-stabddg-K96A-R73A-20260908
\end{verbatim}
Tamarind reported creation at \texttt{2026-09-08 20:03:26} (timestamp as returned by the API)
and subsequently reported completion at \texttt{20:06:59}, after a start at
\texttt{20:06:29}. The recorded queue wait is 183 s, execution interval 30 s, and usage
0.01 weighted hours; the latter is a service accounting unit, not a dollar cost. This run uses Tamarind; the separately staged
100-frame MM-GBSA resample remains an independent Gemini task run by the user.

\paragraph{Reproduction and provenance limits.}
All paths below are relative to the repository root. The named driver is
\texttt{episcaf\_v3/energetics/tamarind/pilot.py}. Its original preparation, validation,
submission, and status commands, together with the retrieval commands, are recorded in
\texttt{energetics/tamarind/README.md} and the manuscript figure/command ledger.
For a compact invocation:
\begin{verbatim}
cd episcaf_v3/energetics/tamarind
python3 pilot.py verify
python3 pilot.py status
# Once the service reports Complete:
python3 pilot.py download
\end{verbatim}
The \texttt{3hfm\_pilot/} directory retains the original request, input manifest, schema,
validation response, normalized submitted request, submission-attempt guard, submission receipt,
and latest status. The completed archive is retained as \texttt{results.zip}, with its checksum in
\texttt{summary.json}. The report script verifies that the returned original PDB matches
the submitted checksum and that both returned mutation strings are unchanged. It also checks
the returned chain selections, sample count, and seed against the frozen request. The service
renumbering log counts 130 chain-Y records, whereas the processed PDB has 129 distinct
\texttt{ATOM} residues; the original input includes a crystal water. The target amino-acid
identities in the processed structure are verified explicitly. The driver prevents accidental duplicate submission of this frozen run.
A new reproduction requires a distinct run directory and service job name, documented in the
README, preserving the original artifacts. Credentials stay outside version control.

A fixed seed and preserved request make the submitted calculation auditable, but do not guarantee
bitwise reproducibility on a hosted service. The log identifies the loaded checkpoint as \texttt{model\_ckpts/stabddg.pt}, but does
not provide its hash or a code/environment version. Service preprocessing is partly documented
by the returned structures and log. The upstream repository describes sequential fine-tuning
on folding stability data and SKEMPI binding data
(\url{https://github.com/LDeng0205/StaB-ddG}); whether these specific mutations were used
to train the deployed checkpoint remains unverified. These are required provenance checks before calling this independent validation.
The upstream documentation defines output in kcal/mol with negative values stabilizing binding,
consistent with our experimental convention; the service log shows that its raw \texttt{pred\_1}
values are copied into \texttt{Prediction}. The two returned values and their signed errors are
reported in Table~\ref{tab:tamarind_pilot}. StaB-ddG orders K96A above R73A, but underestimates
K96A by 4.73 kcal/mol and predicts the opposite sign for R73A (error $+0.81$ kcal/mol).
This is useful separation of the selected pair, not evidence that the model is quantitatively
calibrated or superior across the epitope. Twenty Monte Carlo samples provide one averaged output
per mutant here; no per-sample variance or confidence interval is returned.

\begin{table}[htbp]
\centering
\begin{tabular}{lrrr}
\toprule
Mutation & Experiment & StaB-ddG & Prediction $-$ experiment \\
\midrule
K96A & +6.49 & +1.76 & -4.73 \\
R73A & -0.33 & +0.48 & +0.81 \\
\bottomrule
\end{tabular}

\caption{Two separate antigen alanine substitutions in 3HFM. All entries are in kcal/mol;
positive $\Delta\Delta G$ weakens binding. Generated directly from the downloaded output and
stored experimental CSV by \texttt{energetics/tamarind/report\_pilot.py}.}
\label{tab:tamarind_pilot}
\end{table}

\begin{figure}[htbp]
\centering
\figorbox{tamarind_vs_simulation_3hfm.png}{1.0}
\caption{The learned predictor and simulation signals against experimental antigen alanine
$\Delta\Delta G$. Blue and orange identify K96A and R73A in every panel. The learned model has
only these two selected predictions; the simulation panels include all thirteen measured
residues. The identity line applies only to predicted mutational $\Delta\Delta G$. Simulation
energies are sign-flipped so more favorable contributions point upward, with their own units
and independent vertical scales. Rank correlations use all thirteen measured residues; no
correlation is reported for the two-point ML diagnostic. Generated by
\texttt{energetics/tamarind/plot\_comparison.py}; plotted values are saved as a CSV.}
\label{fig:tamarind_simulation}
\end{figure}

A broader comparison would use all thirteen stored antigen mutations and additional complexes
with an explicit training-overlap audit. No classification performance or category threshold is
estimated from this selected two-point pilot.
